\originalchapter{冷碰撞与量子气体}
\originalsection{散射长度与可调相互作用}
低温时相对运动波长远大于作用范围, 高分波受离心势抑制. $s$波远场径向解$u(r)\propto\sin(kr+\delta_0)$, 定义$a=-\lim_{k\to0}\tan\delta_0/k$. 有效程展开为$k\cot\delta_0=-1/a+r_ek^2/2+\cdots$, 散射幅度$f_0=1/(k\cot\delta_0-\ii k)\simeq-a/(1+\ii ka)$. 可分辨粒子截面$4\pi a^2/(1+k^2a^2)$; 同内部态玻色子在低能下为$8\pi a^2$, 同内部态费米子禁戒$s$波.

两通道闭通道束缚态接近开通道阈值时, $a(B)=a_{\rm bg}[1-\Delta_B/(B-B_0)]$. 这是Feshbach共振的低能参数化. $|a|\to\infty$并不表示有限动量截面无限, 因幺正性把$s$波截面限制为$4\pi/k^2$. 共振附近还须考虑有效程、非弹性损失与三体过程.

\originalsection{蒸发冷却与玻色凝聚}
蒸发冷却移除能量高于$U_c=\eta k_BT$的尾部原子, 再由弹性碰撞重新热化. 去掉高于平均能量的粒子会降低剩余温度, 但原子数也下降. 能否持续增加相空间密度取决于弹性碰撞率能否保持, 背景损失与三体损失是否足够低; 降低阱深不自动等同高效蒸发.

三维均匀理想Bose气的激发态密度为$n_{\rm ex}=\lambda_T^{-3}g_{3/2}(z)$, $\lambda_T=h/\sqrt{2\pi mk_BT}$,$g_s(z)=\sum_{l\ge1}z^l/l^s$. 随$\mu\to0^-$,$z\to1$时$n_{\rm ex}$达到$\zeta(3/2)/\lambda_T^3$, 多余粒子进入最低模式. 故$n\lambda_T^3=\zeta(3/2)$是理想均匀气的临界条件; 谐振阱态密度不同, 临界温度改为$k_BT_c=\hbar\bar\omega[N/\zeta(3)]^{1/3}$. 这些结论的系综与几何必须注明.

\originalsection{相互作用、声子与光晶格}
稀薄弱相互作用Bose气用$g=4\pi\hbar^2a/m$的接触耦合. 能量泛函$E=\int[\hbar^2|\nabla\psi|^2/(2m)+V|\psi|^2+g|\psi|^4/2]\dd^3r$在粒子数约束下变分, 得Gross--Pitaevskii方程$\ii\hbar\partial_t\psi=[-\hbar^2\nabla^2/(2m)+V+g|\psi|^2]\psi$. 令均匀解附近扰动含正负频率两个分量, 线性化二乘二矩阵的行列式给
\[\epsilon_k=\sqrt{E_k(E_k+2gn)},\qquad E_k=\frac{\hbar^2k^2}{2m}.\]
小$k$为声子$\epsilon_k\simeq\hbar ck$, $c=\sqrt{gn/m}$; 大$k$回到粒子. 负$a$的均匀模型会有长波不稳定, 不能任意用同一实声速公式.

深光晶格的最低能带可用Bose--Hubbard模型$H=-J\sum_{\langle ij\rangle}(b_i^\dagger b_j+\hc)+(U/2)\sum_i n_i(n_i-1)-\mu\sum_in_i$. $J$促进离域与相位相干, $U$抑制占据数涨落. 整数填充且$U/J$大时有Mott区; 非整数填充和有限温度会改变相图. 本卷讨论模型来源与两极限, 不声称给出全部临界指数的推导.

\originalsection{Fermi气与BEC--BCS跨越}
两自旋分量的均匀理想气有$E_F=\hbar^2(3\pi^2n)^{2/3}/(2m)$, 这里$n$为两分量总密度. Pauli占据造成Fermi压力, 单一内部态的$s$波禁戒不表示两分量也无相互作用. $a<0$弱吸引的一侧形成重叠Cooper对,$a>0$且$ka\ll1$可形成结合能$\hbar^2/(ma^2)$的浅二聚体. 跨越区参数为$1/(k_Fa)$, 单位极限为零.

BCS平均场的准粒子谱$E_k=\sqrt{(\epsilon_k-\mu)^2+|\Delta|^2}$来自对$\ket{0},\ket{\uparrow\downarrow}$配对基的Bogoliubov对角化. 它说明配对能隙, 但强耦合单位区缺乏小参数, 平均场定量值需实验或更精确多体计算. 这一区间与热统卷中的理想气体相衔接, 不把理想气的表达式当作强相互作用结论.

\originalsection{蒸发效率与runaway的定量条件}
8.422第18讲用能量与碰撞率账讨论持续蒸发. 先取三维经典谐振阱, 频率固定, 截断参数$\eta=U_c/(k_BT)\gg1$. 未截断平均能量$E=3Nk_BT$, 每个逃走原子平均携走$\epsilon_{\rm loss}k_BT=(\eta+\kappa)k_BT$, 深阱极限$\kappa$近1; 有限$\eta$时截断态密度修正这两个系数. 只计有用蒸发, $dE=\epsilon_{\rm loss}k_BT\,dN$, 得
\[\alpha_{\rm ev}\equiv\frac{\dd\ln T}{\dd\ln N}
 =\frac{\epsilon_{\rm loss}-3}{3},\qquad T\propto N^{\alpha_{\rm ev}}.\]
谐振阱中心相空间密度$\mathcal D=N(\hbar\bar\omega/k_BT)^3$, 效率定义
\[\gamma_{\rm ev}=-\frac{\dd\ln\mathcal D}{\dd\ln N}
 =3\alpha_{\rm ev}-1=\epsilon_{\rm loss}-4.\]
高$\eta$有较高每原子效率, 但尾部概率约$\e^{-\eta}$, 蒸发会很慢, 不是越高越好. 热化时间约几次弹性碰撞时间, 需要移除尾部的节奏慢于重新热化.

固定低能$s$波截面时, 中心密度$n_0\propto N\bar\omega^3T^{-3/2}$, 平均速度$\bar v\propto T^{1/2}$, 所以
\[\Gamma_{\rm el}\sim n_0\sigma\bar v\propto N\bar\omega^3/T.\]
固定$\bar\omega$下$\Gamma_{\rm el}\propto N^{1-\alpha_{\rm ev}}$. 原子数递减时碰撞率反而增大须$\alpha_{\rm ev}>1$, 即$\epsilon_{\rm loss}>6$. 这是这一模型的runaway条件, 不是仅有温度降低就称runaway.

若蒸发过程中$\bar\omega\propto T^\nu$, 则
\[\gamma_{\rm ev}=3(1-\nu)\alpha_{\rm ev}-1,\qquad
 \Gamma_{\rm el}\propto N^{\,1+(3\nu-1)\alpha_{\rm ev}}.\]
降低光阱功率时常有$U_c\propto P$, $\bar\omega\propto\sqrt P$, 固定$\eta$令$\nu=1/2$, 此时碰撞率指数为$1+\alpha_{\rm ev}/2>0$, 通常不会runaway. 保持额外约束的磁阱或交叉光阱可改变这个情况.
\begin{example}取$\eta=8,\kappa\simeq1$, 则$\alpha_{\rm ev}\simeq2$. 固定频率阱的效率约5, 弹性碰撞率随$N^{-1}$增长; 同样参数而$\nu=1/2$时效率约2, 碰撞率随$N^2$下降. 两个实验即便同样降温, 蒸发动力学也很不同.\end{example}
背景碰撞把平均能量附近的原子随机移除, 会损失$N$而少改$T$, 降低效率. 三体损失在高密度处尤其严重, 并可选择性损失低势能原子而加热. 要持续提高$\mathcal D$, 有用蒸发/热化率须压过背景与非弹性损失, 且散射截面尚未因幺正限、温度或窄共振改变. 上述幂律的近似条件在跨入退化区后也须重新检查.

\originalsection{愈合长度与Thomas--Fermi剖面}
以下取排斥$a>0$的稀薄三维气. 均匀凝聚$\mu=gn_0$, 在边界扰动处令$\psi=\sqrt{n_0}f(x)$, 稳态GP式变为
\[-\xi^2 f''+(f^2-1)f=0,\qquad
 \xi=\frac{\hbar}{\sqrt{2mgn_0}}.\]
半空间硬墙$f(0)=0,f(\infty)=1$的解为$f=\tanh[x/(\sqrt2\xi)]$. 因而密度不是到边界瞬间截断, 动能与相互作用在$\xi$尺度平衡. 均匀Bogoliubov谱的交叉也在$k\xi$约1: 长波为声子, 短波为粒子.

阱内相互作用主导且密度慢变时忽略梯度能, 得Thomas--Fermi(TF)密度
\[n(\mathbf r)=\frac{\mu-V(\mathbf r)}g\quad(V<\mu),\qquad n=0\quad(V\ge\mu).\]
三维各向同性谐振阱$V=m\omega^2r^2/2$有$R^2=2\mu/(m\omega^2)$,
\[N=\frac{8\pi\mu R^3}{15g},\quad
 \mu=\frac{\hbar\omega}{2}\left(\frac{15Na}{a_{\rm ho}}\right)^{2/5},\quad
 R=a_{\rm ho}\left(\frac{15Na}{a_{\rm ho}}\right)^{1/5},
 \quad a_{\rm ho}=\sqrt{\frac{\hbar}{m\omega}}.\]
积分中的$8\pi/15$来自抛物线剖面, 不是Gaussian宽度. TF要求$Na/a_{\rm ho}\gg1$, 稀薄气同时要求中心$n_0a^3\ll1$. 阱边缘仍由愈合层平滑, 在那里删掉梯度项不再可靠.

\originalsection{Bogoliubov振幅与零温耗尽}
令场算符$\hat\psi=\sqrt{n_0}+\delta\hat\psi$, 保留二阶涨落. 对$(a_{\mathbf k},a_{-\mathbf k}^\dagger)$的运动矩阵为
\[\begin{pmatrix}E_k+gn_0&gn_0\\-gn_0&-E_k-gn_0\end{pmatrix},\qquad
 \epsilon_k^2=E_k(E_k+2gn_0).\]
以$a_{\mathbf k}=u_kb_{\mathbf k}-v_kb_{-\mathbf k}^\dagger$作正则Bogoliubov变换, 保持玻色对易子要求$u_k^2-v_k^2=1$, 并得
\[u_k^2=\frac12\left(\frac{E_k+gn_0}{\epsilon_k}+1\right),\qquad
 v_k^2=\frac12\left(\frac{E_k+gn_0}{\epsilon_k}-1\right).\]
即使准粒子真空$b_{\mathbf k}\ket0=0$, 原子模仍有$\langle a_{\mathbf k}^\dagger a_{\mathbf k}\rangle=v_k^2$. 它是相互作用产生的量子耗尽, 与有限温热激发不同.

三维连续积分令$q=k\xi$, 有
\[\frac{n_{\rm dep}}{n_0}
 =\frac1{2\pi^2n_0\xi^3}\int_0^\infty
 \frac{q^2}{2}\left[\frac{q^2+1}{q\sqrt{q^2+2}}-1\right]\dd q
 =\frac8{3\sqrt\pi}\sqrt{n_0a^3}.\]
无量纲积分为$1/(3\sqrt2)$. 小耗尽与$n_0a^3\ll1$一致, 这才自洽地支持凝聚背景展开. 一维、二维红外行为、有限温度及强相互作用不能直接使用此三维零温式.

\originalsection{Bose--Hubbard的平均场Mott边界}
把最低能带Wannier函数$w_i$代入接触相互作用, 得$U=g\int|w_i|^4\dd^3r$, 相邻单粒子矩阵元给$J$. 固定整数填充$n$, $J=0$时每站点能$\mathcal E_n=Un(n-1)/2-\mu n$, 稳定区$U(n-1)<\mu<Un$. 粒子和空穴能隙分别$\Delta_p=Un-\mu,\Delta_h=\mu-U(n-1)$.

取均匀$\phi=\langle b_i\rangle$实, 协调数$z$, 分解$b_i^\dagger b_j\simeq\phi(b_i^\dagger+b_j)-\phi^2$. 单站点Hamiltonian为
\[h_{\rm MF}=Un(n-1)/2-\mu n-zJ\phi(b^\dagger+b)+zJ\phi^2.\]
在$|n\rangle$上二阶微扰, $\phi^2$系数为
\[zJ-(zJ)^2\left(\frac{n+1}{Un-\mu}
                         +\frac n{\mu-U(n-1)}\right).\]
系数变号时$\phi=0$失稳. 写$x=\mu/U$, 边界为
\[\frac{zJ_c}U=\frac{(n-x)(x-n+1)}{x+1},\qquad n-1<x<n.\]
顶点最大值$(zJ/U)_{\rm tip}=(\sqrt{n+1}-\sqrt n)^2$. 例如$n=1$顶点约0.1716, 这是平均场值而非任何维度的精确实验临界比. Mott区有整数占据和粒子/空穴隙; 非整数填充的均匀系统不落在此$J=0$绝缘区. 强涨落和低维会改变边界, 但这个计算给出了竞争怎样变成可画的相图.

\originalsection{Cooper不稳定性与BCS波函数}
一个Fermi海上方的两异自旋粒子, 在$0<\xi_k<\hbar\omega_c$壳内有常数吸引$-g/V$($g>0$). 对总动量0的配对振幅$a_k$, 束缚能$E_B>0$满足
\[(2\xi_k+E_B)a_k=\frac gV\sum_{k'}a_{k'},\qquad
 1\simeq\frac{g\nu_0}{2}\log\frac{2\hbar\omega_c+E_B}{E_B}.\]
这里$\nu_0$是每体积每自旋的Fermi面态密度. 弱吸引给$E_B\simeq2\hbar\omega_c\exp[-2/(g\nu_0)]$, 任意吸引都产生解. 对数来自Fermi面附近的连续配对态, 不是孤立两个自由粒子在三维总能束缚的判据; Pauli海和壳截断是模型的关键条件.

许多重叠对不能各自当作可辨认二聚体. BCS变分态写为
\[\ket{\mathrm{BCS}}=\prod_{\mathbf k}
 (u_k+v_kc_{\mathbf k\uparrow}^\dagger
                  c_{-\mathbf k\downarrow}^\dagger)\ket0,\qquad
 |u_k|^2+|v_k|^2=1.\]
每对模式在空与双占据之间相干叠加, Pauli限制自动包括在费米算符中. 对平均场能量或两维配对块变分, 可选$\Delta$实正并得到
\[E_k=\sqrt{\xi_k^2+\Delta^2},\quad
 u_k^2=\tfrac12(1+\xi_k/E_k),\quad
 v_k^2=\tfrac12(1-\xi_k/E_k),\quad
 u_kv_k=\Delta/(2E_k),\]
\[\Delta=\frac gV\sum_ku_kv_k,\qquad N=2\sum_kv_k^2.\]
恒定态密度壳模型的零温隙$\Delta\simeq2\hbar\omega_c\e^{-1/(g\nu_0)}$, 指数与单Cooper对问题不同. 准粒子激发需破坏配对, 产生能隙, 这联系原讲义中的超流而不是仅给一条谱公式.

冷原子的接触吸引须用散射长度消除紫外截止. 平均场的有限形式为
\[\frac{m}{4\pi\hbar^2a}=\int\frac{\dd^3k}{(2\pi)^3}
 \left(\frac1{2\epsilon_k^0}-\frac1{2E_k}\right),\qquad
 n=\int\frac{\dd^3k}{(2\pi)^3}(1-\xi_k/E_k),\quad
 \epsilon_k^0=\hbar^2k^2/(2m).\]
同时求$\mu,\Delta$可连接BCS与分子BEC两端. $a>0$浅二聚体的$\hbar^2/(ma^2)$与$1/(k_Fa)$控制跨越; 单位极限没有小参数, 上述平均场只给定性组织, 精确能量、临界温度和涨落需更完整多体证据.
